%% GENERATED by generate_papers.py -- edit the bank_*.py files, not this file.
\documentclass[11pt,a4paper]{article}
\usepackage[a4paper, top=20mm, bottom=20mm, left=18mm, right=18mm, headheight=14pt]{geometry}
\usepackage[T1]{fontenc}
\usepackage{lmodern}
\usepackage{amsmath, amssymb, mathtools}
\usepackage{braket}
\usepackage{array, tabularx, booktabs, multicol}
\usepackage{enumitem}
\usepackage[dvipsnames]{xcolor}
\usepackage[most]{tcolorbox}
\usepackage{needspace}
\usepackage{fancyhdr, lastpage}
\usepackage{pifont}
\usepackage{tikz}
\usepackage{graphicx}
\usepackage{eso-pic}
% SILICO FELLER watermark: official logo, centred, ~8 % opacity, behind all content, on every page
\AddToShipoutPictureBG{\AtPageCenter{\makebox(0,0){\tikz\node[opacity=0.08, inner sep=0pt]{\includegraphics[width=0.72\paperwidth]{../../../assets/branding/silico_feller_logo.png}};}}}
\usepackage[hidelinks]{hyperref}
\newcolumntype{L}{>{\raggedright\arraybackslash}X}
\definecolor{brand}{HTML}{1F4E79}
\definecolor{okgreen}{HTML}{1B7F3B}
\setlength{\parindent}{0pt}
\setlength{\parskip}{3pt}
\newcommand{\optlabel}[1]{\textbf{(#1)}}
\newcommand{\correctlabel}[1]{\textcolor{okgreen}{\textbf{(#1)}\,\ding{51}}}
\newenvironment{qitem}[1]{\begin{list}{}{\setlength{\leftmargin}{3.3em}\setlength{\labelwidth}{3.1em}%
  \setlength{\labelsep}{0.2em}\setlength{\topsep}{3pt}\setlength{\parsep}{3pt}}\item[\textbf{Q.#1}\hfill]}{\end{list}}
\newcommand{\qsep}{\par\vspace{2pt}{\color{black!15}\hrule}\vspace{4pt}}
\newcommand{\qmeta}[2]{{\footnotesize\color{black!60}\textsc{#1} \;$\cdot$\; #2}\par}
\newtcolorbox{passagebox}[1]{enhanced, breakable, colback=brand!4, colframe=brand!70, fonttitle=\bfseries,
  title={#1}, left=6pt, right=6pt, top=4pt, bottom=4pt}
\newtcolorbox{sectionbanner}{enhanced, colback=brand, colframe=brand, coltext=white, fontupper=\large\bfseries,
  left=8pt, right=8pt, top=6pt, bottom=6pt, sharp corners}
\newtcolorbox{infobox}[1]{enhanced, breakable, colback=black!3, colframe=black!45, fonttitle=\bfseries,
  title={#1}, left=6pt, right=6pt, top=4pt, bottom=4pt}
\newtcolorbox{solbox}[1]{enhanced, breakable, colback=okgreen!3, colframe=okgreen!60!black, fonttitle=\bfseries,
  title={Answer: (#1) \quad Detailed solution}, left=6pt, right=6pt, top=3pt, bottom=4pt, before skip=4pt, after skip=6pt}
\allowdisplaybreaks
\setlength{\emergencystretch}{3em}

\begin{document}
\pagestyle{fancy}\fancyhf{}
\fancyhead[L]{\small\color{black!70}Quantum Noise Simulation --- Assessment}
\fancyhead[R]{\small\color{black!70}\textbf{Question Paper}}
\fancyfoot[C]{\small Page \thepage\ of \pageref*{LastPage}}
\renewcommand{\headrulewidth}{0.4pt}

\thispagestyle{empty}
\begin{center}
{\color{brand}\rule{\linewidth}{1.2pt}}\\[6pt]
{\Huge\bfseries\color{brand} Quantum Noise Simulation}\\[4pt]
{\Large Assessment of noise models, channels, simulation and mitigation}\\[4pt]
{\large\bfseries QUESTION PAPER}\\[2pt]
{\color{brand}\rule{\linewidth}{1.2pt}}\\[1pt]
\end{center}
\vspace{-2pt}
\begin{tabularx}{\linewidth}{@{}lXlr@{}}
\textbf{Time allowed:} & 3 hours & \textbf{Maximum marks:} & 240\\
\textbf{Questions:} & 80 (4 sections $\times$ 20) & \textbf{Answer key:} & separate Answer Book
\end{tabularx}
\vspace{6pt}

\textbf{Name of the candidate:} \hrulefill\quad \textbf{Roll No.:} \rule{35mm}{0.4pt}

\vspace{6pt}
\begin{infobox}{General instructions}
\begin{enumerate}[leftmargin=1.4em, itemsep=1pt, topsep=1pt]
\item This question paper has \textbf{four sections} of \textbf{20 questions} each:
\textbf{A}~Mathematical and Quantum Foundations, \textbf{B}~Quantum Noise Fundamentals, \textbf{C}~Noise Models and Simulation, \textbf{D}~Advanced Noise Analysis and Implementation.
\item Every question is of the \textbf{single-correct} type: four options, exactly one correct. Question formats used:
single-correct MCQ, matching lists (List-I/List-II), Assertion--Reason, Statement~I/II and paragraph-based questions.
\item For every question choose \textbf{ONE} option, (A), (B), (C) or (D), and record it against the question number in the answer script provided by the invigilator.
\item Calculators are \textbf{not} allowed. Rough work may be done on separate sheets.
\item Unless a question says otherwise, use the \textbf{conventions} printed below; they apply to every section.
\item The detailed solutions are in the separate \emph{Answer Book}; do not consult it before finishing the test.
\end{enumerate}
\end{infobox}
\begin{infobox}{Marking scheme (every question)}
Each question has \textbf{four} options (A), (B), (C) and (D). \textbf{ONLY ONE} of these four options is correct.\\[2pt]
\begin{tabular}{@{}ll@{}}
\emph{Full marks}: & $+3$ if ONLY the correct option is chosen;\\
\emph{Zero marks}: & $\phantom{+}0$ if none of the options is chosen (question unanswered);\\
\emph{Negative marks}: & $-1$ in all other cases (wrong option, or more than one option chosen).
\end{tabular}
\end{infobox}

\begin{infobox}{Conventions and notation used throughout the paper}
\footnotesize
\begin{itemize}[leftmargin=1.2em, itemsep=1pt, topsep=2pt]
\item Qubit 0 is written first: $\ket{q_0q_1\dots}$; two-qubit matrices use the basis order $\ket{00}, \ket{01}, \ket{10}, \ket{11}$. $X, Y, Z$ are the Pauli matrices,
      $H$ the Hadamard gate, $R_P(\theta) = e^{-i\theta P/2}$; CNOT flips the target when the control is 1.
\item A qubit state is $\rho = \frac12(I + r_xX + r_yY + r_zZ)$ with Bloch vector $\vec r$, $|\vec r| \le 1$; purity $\operatorname{Tr}\rho^2$. State fidelity with a pure target:
      $F = \braket{\psi|\rho|\psi}$. Measurement probabilities are the diagonal of $\rho$; $\langle Z\rangle = P(0) - P(1)$.
\item A channel is $\mathcal E(\rho) = \sum_kK_k\rho K_k^\dagger$ with $\sum_kK_k^\dagger K_k = I$. Bit flip: $\sqrt{1-p}\,I, \sqrt p\,X$; phase flip: $\sqrt{1-p}\,I, \sqrt p\,Z$;
      bit-phase flip: $\sqrt{1-p}\,I, \sqrt p\,Y$; depolarizing (CUDA-Q): $\sqrt{1-p}\,I, \sqrt{p/3}\,X, \sqrt{p/3}\,Y, \sqrt{p/3}\,Z$ (two-qubit: $\sqrt{1-p}\,II$ and the 15
      other Paulis with $p/15$); amplitude damping: $\operatorname{diag}(1, \sqrt{1-\gamma})$, $\sqrt\gamma\ket0\!\bra1$; phase damping: $\operatorname{diag}(1, \sqrt{1-\lambda})$, $\operatorname{diag}(0, \sqrt\lambda)$.
\item $T_1$: $\rho_{11}(t) = \rho_{11}(0)e^{-t/T_1}$; $T_2$: $|\rho_{01}(t)| = |\rho_{01}(0)|e^{-t/T_2}$; $\frac1{T_2} = \frac1{2T_1} + \frac1{T_\varphi}$. $\ln$ is the natural logarithm.
\item \textbf{Notebook} (\texttt{Noise\_Simulation.ipynb}, defaults): 15-qubit GHZ circuit ($H$ on qubit 0, CNOT chain), the chosen channel with probability $p$ after
      every \texttt{h}, \texttt{x}, \texttt{ry}, \texttt{rz} gate and a two-qubit channel after every CNOT; no measurement noise; 10\,000 shots, seed 42.
      \textbf{Previously recorded run}: CUDA-Q 0.16.0, CPU target \texttt{stim}, no GPU, depolarizing $p = 0.01$: population $0.8936$, $X$-basis parity $0.6910$,
      $(1-p)^G = 0.8601$. No physical quantum processor was used.
\end{itemize}
\end{infobox}
\newpage

\begin{sectionbanner}SECTION A \enspace$\cdot$\enspace Mathematical and Quantum Foundations \hfill Q.1--20\end{sectionbanner}

{\small \textbf{Syllabus:} probability and error statistics, complex numbers and phases, matrices and linear operators, density matrices, pure and mixed states, purity, eigenvalues, partial trace, the Bloch sphere and fidelity of states.\\ \textbf{Formats:} Single-correct MCQ (Q.1--14); Matching lists (Q.15--16); Assertion--Reason (Q.17); Statement I / II (Q.18); Paragraph based (Q.19--20). Each question: $+3$ / $0$ / $-1$.}
\par\medskip

\needspace{7\baselineskip}
\begin{qitem}{1}
The density matrix of the pure state $\ket\psi = \frac{1}{\sqrt2}(\ket0 + i\ket1)$ is
\par\smallskip\noindent\begin{tabularx}{\linewidth}{@{}*{4}{L}@{}}
\optlabel{A}\,$\frac12\begin{pmatrix}1 & 0\\ 0 & 1\end{pmatrix}$ & \optlabel{B}\,$\frac12\begin{pmatrix}1 & 1\\ 1 & 1\end{pmatrix}$ & \optlabel{C}\,$\frac12\begin{pmatrix}1 & -i\\ i & 1\end{pmatrix}$ & \optlabel{D}\,$\frac12\begin{pmatrix}1 & i\\ -i & 1\end{pmatrix}$
\end{tabularx}
\end{qitem}
\qsep

\needspace{7\baselineskip}
\begin{qitem}{2}
The purity $\operatorname{Tr}\rho^2$ of $\rho = \begin{pmatrix}0.7 & 0.3\\ 0.3 & 0.3\end{pmatrix}$ is
\par\smallskip\noindent\begin{tabularx}{\linewidth}{@{}*{4}{L}@{}}
\optlabel{A}\,$0.58$ & \optlabel{B}\,$1$ & \optlabel{C}\,$0.76$ & \optlabel{D}\,$0.52$
\end{tabularx}
\end{qitem}
\qsep

\needspace{7\baselineskip}
\begin{qitem}{3}
The Bloch vector $(r_x, r_y, r_z)$ of $\rho = \begin{pmatrix}0.8 & 0.2 - 0.1i\\ 0.2 + 0.1i & 0.2\end{pmatrix}$ is
\par\smallskip\noindent\begin{tabularx}{\linewidth}{@{}*{2}{L}@{}}
\optlabel{A}\,$(0.2,\ 0.1,\ 0.6)$ & \optlabel{B}\,$(0.4,\ 0.2,\ 0.6)$ \\[3pt]
\optlabel{C}\,$(0.4,\ -0.2,\ 0.6)$ & \optlabel{D}\,$(0.4,\ 0.2,\ 0.8)$
\end{tabularx}
\end{qitem}
\qsep

\needspace{7\baselineskip}
\begin{qitem}{4}
Which of the following matrices is \emph{not} a valid density matrix?
\par\smallskip\noindent\begin{tabularx}{\linewidth}{@{}*{4}{L}@{}}
\optlabel{A}\,$\begin{pmatrix}0.5 & 0.5\\ 0.5 & 0.5\end{pmatrix}$ & \optlabel{B}\,$\begin{pmatrix}0.5 & -0.3i\\ 0.3i & 0.5\end{pmatrix}$ & \optlabel{C}\,$\begin{pmatrix}0.9 & 0\\ 0 & 0.1\end{pmatrix}$ & \optlabel{D}\,$\begin{pmatrix}0.6 & 0.5\\ 0.5 & 0.4\end{pmatrix}$
\end{tabularx}
\end{qitem}
\qsep

\needspace{7\baselineskip}
\begin{qitem}{5}
The operators $K_0 = \begin{pmatrix}1 & 0\\ 0 & 0.8\end{pmatrix}$ and $K_1 = \begin{pmatrix}0 & b\\ 0 & 0\end{pmatrix}$ ($b \ge 0$) form a
trace-preserving set of Kraus operators exactly when $b$ equals
\par\smallskip\noindent\begin{tabularx}{\linewidth}{@{}*{4}{L}@{}}
\optlabel{A}\,$0.6$ & \optlabel{B}\,$0.36$ & \optlabel{C}\,$0.8$ & \optlabel{D}\,$0.4$
\end{tabularx}
\end{qitem}
\qsep

\needspace{7\baselineskip}
\begin{qitem}{6}
A qubit is prepared in $\frac{1}{\sqrt2}(\ket0 + e^{i\varphi}\ket1)$ where the phase $\varphi$ is random and uniformly distributed in
$[-\frac\pi2, \frac\pi2]$. The average value of $\langle X\rangle$ is
\par\smallskip\noindent\begin{tabularx}{\linewidth}{@{}*{4}{L}@{}}
\optlabel{A}\,$\pi/4 \approx 0.79$ & \optlabel{B}\,$1/2$ & \optlabel{C}\,$2/\pi \approx 0.64$ & \optlabel{D}\,$0$
\end{tabularx}
\end{qitem}
\qsep

\needspace{7\baselineskip}
\begin{qitem}{7}
A circuit has 20 gates; each fails independently with probability $p = 0.01$. The probability that \emph{at least one} gate fails is
\par\smallskip\noindent\begin{tabularx}{\linewidth}{@{}*{4}{L}@{}}
\optlabel{A}\,$0.182$ & \optlabel{B}\,$0.818$ & \optlabel{C}\,$0.200$ & \optlabel{D}\,$0.165$
\end{tabularx}
\end{qitem}
\qsep

\needspace{7\baselineskip}
\begin{qitem}{8}
An outcome has true probability $0.9$. It is estimated from $10\,000$ shots. The standard error of the estimate is
\par\smallskip\noindent\begin{tabularx}{\linewidth}{@{}*{4}{L}@{}}
\optlabel{A}\,$0.03$ & \optlabel{B}\,$0.003$ & \optlabel{C}\,$0.0009$ & \optlabel{D}\,$0.009$
\end{tabularx}
\end{qitem}
\qsep

\needspace{7\baselineskip}
\begin{qitem}{9}
An $X$ gate is implemented as $R_X(\pi + \varepsilon)$ instead of $R_X(\pi)$, with $\varepsilon = 0.1$ rad. Applied to $\ket0$, the
probability of measuring $0$ (an error) is closest to
\par\smallskip\noindent\begin{tabularx}{\linewidth}{@{}*{4}{L}@{}}
\optlabel{A}\,$0.0025$ & \optlabel{B}\,$0.01$ & \optlabel{C}\,$0.1$ & \optlabel{D}\,$0.05$
\end{tabularx}
\end{qitem}
\qsep

\needspace{7\baselineskip}
\begin{qitem}{10}
The fidelity $F = \braket{+|\rho|+}$ between the target $\ket+$ and the mixed state $\rho = 0.8\ket+\!\bra+ + 0.2\,\frac I2$ is
\par\smallskip\noindent\begin{tabularx}{\linewidth}{@{}*{4}{L}@{}}
\optlabel{A}\,$0.6$ & \optlabel{B}\,$0.9$ & \optlabel{C}\,$0.8$ & \optlabel{D}\,$0.95$
\end{tabularx}
\end{qitem}
\qsep

\needspace{7\baselineskip}
\begin{qitem}{11}
For the two-qubit state $\ket\psi = \sqrt{0.8}\ket{00} + \sqrt{0.2}\ket{11}$, the purity $\operatorname{Tr}\rho_A^2$ of the reduced state of qubit A is
\par\smallskip\noindent\begin{tabularx}{\linewidth}{@{}*{4}{L}@{}}
\optlabel{A}\,$0.8$ & \optlabel{B}\,$0.5$ & \optlabel{C}\,$0.68$ & \optlabel{D}\,$1$
\end{tabularx}
\end{qitem}
\qsep

\needspace{7\baselineskip}
\begin{qitem}{12}
The larger eigenvalue of $\rho = \begin{pmatrix}0.7 & 0.3\\ 0.3 & 0.3\end{pmatrix}$ is closest to
\par\smallskip\noindent\begin{tabularx}{\linewidth}{@{}*{4}{L}@{}}
\optlabel{A}\,$0.700$ & \optlabel{B}\,$0.861$ & \optlabel{C}\,$1.000$ & \optlabel{D}\,$0.760$
\end{tabularx}
\end{qitem}
\qsep

\needspace{7\baselineskip}
\begin{qitem}{13}
Two independent qubits are in the states $\rho_A = \operatorname{diag}(0.9, 0.1)$ and $\rho_B = \operatorname{diag}(0.8, 0.2)$. In the
state $\rho_A\otimes\rho_B$, the probability that both qubits give the \emph{same} outcome is
\par\smallskip\noindent\begin{tabularx}{\linewidth}{@{}*{4}{L}@{}}
\optlabel{A}\,$0.74$ & \optlabel{B}\,$0.50$ & \optlabel{C}\,$0.26$ & \optlabel{D}\,$0.72$
\end{tabularx}
\end{qitem}
\qsep

\needspace{7\baselineskip}
\begin{qitem}{14}
A qubit with Bloch vector $(r_x, r_y, r_z)$ is acted on by $\rho \mapsto X\rho X$. Its new Bloch vector is
\par\smallskip\noindent\begin{tabularx}{\linewidth}{@{}*{4}{L}@{}}
\optlabel{A}\,$(-r_x,\ -r_y,\ -r_z)$ & \optlabel{B}\,$(-r_x,\ r_y,\ r_z)$ & \optlabel{C}\,$(r_x,\ -r_y,\ -r_z)$ & \optlabel{D}\,$(r_x,\ r_y,\ -r_z)$
\end{tabularx}
\end{qitem}
\qsep

\needspace{7\baselineskip}
\begin{qitem}{15}
Match each state (List-I) with its Bloch vector (List-II).
\par\smallskip\noindent{\renewcommand{\arraystretch}{1.35}%
\begin{tabularx}{\linewidth}{@{}c>{\raggedright\arraybackslash}p{0.33\linewidth}@{\hspace{1.2em}}cL@{}}
\toprule
\multicolumn{2}{@{}l}{\textbf{List-I}} & \multicolumn{2}{l@{}}{\textbf{List-II}} \\ \midrule
\textbf{(P)} & $\ket1$ & \textbf{(1)} & $(0,\ 0,\ 0.5)$ \\
\textbf{(Q)} & $\frac I2$ & \textbf{(2)} & $(0,\ 0,\ 0)$ \\
\textbf{(R)} & $\frac12\begin{psmallmatrix}1 & -i\\ i & 1\end{psmallmatrix}$ & \textbf{(3)} & $(0,\ 0,\ -1)$ \\
\textbf{(S)} & $\operatorname{diag}(0.75,\ 0.25)$ & \textbf{(4)} & $(0,\ 1,\ 0)$ \\
\bottomrule
\end{tabularx}}
\par\smallskip\noindent\begin{tabularx}{\linewidth}{@{}*{2}{L}@{}}
\optlabel{A}\,P$\to$3, Q$\to$1, R$\to$4, S$\to$2 & \optlabel{B}\,P$\to$3, Q$\to$2, R$\to$1, S$\to$4 \\[3pt]
\optlabel{C}\,P$\to$1, Q$\to$2, R$\to$4, S$\to$3 & \optlabel{D}\,P$\to$3, Q$\to$2, R$\to$4, S$\to$1
\end{tabularx}
\end{qitem}
\qsep

\needspace{7\baselineskip}
\begin{qitem}{16}
Match each matrix (List-I) with its property (List-II).
\par\smallskip\noindent{\renewcommand{\arraystretch}{1.35}%
\begin{tabularx}{\linewidth}{@{}c>{\raggedright\arraybackslash}p{0.33\linewidth}@{\hspace{1.2em}}cL@{}}
\toprule
\multicolumn{2}{@{}l}{\textbf{List-I}} & \multicolumn{2}{l@{}}{\textbf{List-II}} \\ \midrule
\textbf{(P)} & $\begin{psmallmatrix}0 & 1\\ 1 & 0\end{psmallmatrix}$ & \textbf{(1)} & unitary but not Hermitian \\
\textbf{(Q)} & $\frac12\begin{psmallmatrix}1 & 1\\ 1 & 1\end{psmallmatrix}$ & \textbf{(2)} & Hermitian and unitary \\
\textbf{(R)} & $\begin{psmallmatrix}1 & 0\\ 0 & i\end{psmallmatrix}$ & \textbf{(3)} & neither Hermitian nor unitary (a Kraus operator of amplitude damping, up to a factor) \\
\textbf{(S)} & $\begin{psmallmatrix}0 & 1\\ 0 & 0\end{psmallmatrix}$ & \textbf{(4)} & a rank-1 projector and a valid density matrix \\
\bottomrule
\end{tabularx}}
\par\smallskip\noindent\begin{tabularx}{\linewidth}{@{}*{2}{L}@{}}
\optlabel{A}\,P$\to$4, Q$\to$2, R$\to$1, S$\to$3 & \optlabel{B}\,P$\to$2, Q$\to$4, R$\to$3, S$\to$1 \\[3pt]
\optlabel{C}\,P$\to$1, Q$\to$4, R$\to$2, S$\to$3 & \optlabel{D}\,P$\to$2, Q$\to$4, R$\to$1, S$\to$3
\end{tabularx}
\end{qitem}
\qsep

\needspace{7\baselineskip}
\begin{qitem}{17}
\textbf{Assertion (A):} Every mixed single-qubit state lies strictly inside the Bloch ball ($|\vec r| < 1$).\\[1mm]
\textbf{Reason (R):} For $\rho = \frac12(I + \vec r\cdot\vec\sigma)$ one has $\operatorname{Tr}\rho^2 = \frac{1 + |\vec r|^2}{2}$, which equals 1 only when $|\vec r| = 1$.
\begin{itemize}[leftmargin=2.6em, labelsep=0.4em, itemsep=2pt, topsep=3pt]
\item[\optlabel{A}] Both Assertion (A) and Reason (R) are true, and R is the correct explanation of A
\item[\optlabel{B}] Both A and R are true, but R is \emph{not} the correct explanation of A
\item[\optlabel{C}] A is true, but R is false
\item[\optlabel{D}] A is false, but R is true
\end{itemize}
\end{qitem}
\qsep

\needspace{7\baselineskip}
\begin{qitem}{18}
\textbf{Statement I:} Every $2\times2$ Hermitian matrix with trace 1 is a valid density matrix.\\[1mm]
\textbf{Statement II:} The diagonal entries of a density matrix in the computational basis are the probabilities of the outcomes of a
computational-basis measurement.
\begin{itemize}[leftmargin=2.6em, labelsep=0.4em, itemsep=2pt, topsep=3pt]
\item[\optlabel{A}] Both Statement I and Statement II are true
\item[\optlabel{B}] Both Statement I and Statement II are false
\item[\optlabel{C}] Statement I is true, but Statement II is false
\item[\optlabel{D}] Statement I is false, but Statement II is true
\end{itemize}
\end{qitem}
\qsep

\needspace{20\baselineskip}
\begin{passagebox}{Paragraph for Questions 19 and 20}
A qubit should have been prepared in $\ket+ = \frac{1}{\sqrt2}(\ket0 + \ket1)$. State tomography of the real (noisy)
preparation gives the density matrix
\[ \rho = \begin{pmatrix}0.55 & 0.35\\ 0.35 & 0.45\end{pmatrix}. \]
\end{passagebox}
\needspace{7\baselineskip}
\begin{qitem}{19}
The Bloch vector and the purity of the prepared state are
\par\smallskip\noindent\begin{tabularx}{\linewidth}{@{}*{2}{L}@{}}
\optlabel{A}\,$(0.7,\ 0,\ 0.1)$ and $0.50$ & \optlabel{B}\,$(0.7,\ 0,\ 0.55)$ and $0.85$ \\[3pt]
\optlabel{C}\,$(0.35,\ 0,\ 0.1)$ and $0.75$ & \optlabel{D}\,$(0.7,\ 0,\ 0.1)$ and $0.75$
\end{tabularx}
\end{qitem}
\qsep

\needspace{7\baselineskip}
\begin{qitem}{20}
The fidelity $\braket{+|\rho|+}$ of the preparation, and the probability that a computational-basis measurement gives $1$, are
\par\smallskip\noindent\begin{tabularx}{\linewidth}{@{}*{4}{L}@{}}
\optlabel{A}\,$0.70$ and $0.45$ & \optlabel{B}\,$0.85$ and $0.45$ & \optlabel{C}\,$0.92$ and $0.45$ & \optlabel{D}\,$0.85$ and $0.55$
\end{tabularx}
\end{qitem}
\qsep

\newpage

\begin{sectionbanner}SECTION B \enspace$\cdot$\enspace Quantum Noise Fundamentals \hfill Q.21--40\end{sectionbanner}

{\small \textbf{Syllabus:} decoherence, bit flips, phase flips, bit-phase flips, depolarizing noise, amplitude and phase damping, $T_1$, $T_2$ and pure dephasing, measurement (readout) errors, coherent versus stochastic errors.\\ \textbf{Formats:} Single-correct MCQ (Q.21--34); Matching lists (Q.35--36); Assertion--Reason (Q.37); Statement I / II (Q.38); Paragraph based (Q.39--40). Each question: $+3$ / $0$ / $-1$.}
\par\medskip

\needspace{7\baselineskip}
\begin{qitem}{21}
Which single-qubit state is left \emph{unchanged} by the bit-flip channel $\rho \mapsto (1 - p)\rho + pX\rho X$ for every $p$?
\par\smallskip\noindent\begin{tabularx}{\linewidth}{@{}*{2}{L}@{}}
\optlabel{A}\,$\frac{1}{\sqrt2}(\ket0 + i\ket1)$ & \optlabel{B}\,$\ket1$ \\[3pt]
\optlabel{C}\,$\ket0$ & \optlabel{D}\,$\ket+$
\end{tabularx}
\end{qitem}
\qsep

\needspace{7\baselineskip}
\begin{qitem}{22}
A phase-flip channel with $p = 0.1$ acts on $\ket+$. The state is then measured in the $X$ basis. The probability of the
outcome ``$+$'' is
\par\smallskip\noindent\begin{tabularx}{\linewidth}{@{}*{4}{L}@{}}
\optlabel{A}\,$0.95$ & \optlabel{B}\,$1.0$ & \optlabel{C}\,$0.8$ & \optlabel{D}\,$0.9$
\end{tabularx}
\end{qitem}
\qsep

\needspace{7\baselineskip}
\begin{qitem}{23}
A depolarizing channel $\rho \mapsto (1 - p)\rho + \frac p3(X\rho X + Y\rho Y + Z\rho Z)$ with $p = 0.03$ acts on $\ket0$. The probability of
then measuring $1$ is
\par\smallskip\noindent\begin{tabularx}{\linewidth}{@{}*{4}{L}@{}}
\optlabel{A}\,$0.03$ & \optlabel{B}\,$0.04$ & \optlabel{C}\,$0.02$ & \optlabel{D}\,$0.01$
\end{tabularx}
\end{qitem}
\qsep

\needspace{7\baselineskip}
\begin{qitem}{24}
Amplitude damping with $\gamma = 0.36$ acts on $\ket+$. The resulting density matrix is
\par\smallskip\noindent\begin{tabularx}{\linewidth}{@{}*{4}{L}@{}}
\optlabel{A}\,$\begin{pmatrix}0.68 & 0.40\\ 0.40 & 0.32\end{pmatrix}$ & \optlabel{B}\,$\begin{pmatrix}0.68 & 0.32\\ 0.32 & 0.32\end{pmatrix}$ & \optlabel{C}\,$\begin{pmatrix}0.50 & 0.40\\ 0.40 & 0.50\end{pmatrix}$ & \optlabel{D}\,$\begin{pmatrix}0.68 & 0.50\\ 0.50 & 0.32\end{pmatrix}$
\end{tabularx}
\end{qitem}
\qsep

\needspace{7\baselineskip}
\begin{qitem}{25}
Phase damping with parameter $\lambda$ multiplies the off-diagonal elements by $\sqrt{1 - \lambda}$; a phase-flip channel multiplies them by $1 - 2p$.
The phase-flip probability $p$ that has the same effect as phase damping with $\lambda = 0.36$ is
\par\smallskip\noindent\begin{tabularx}{\linewidth}{@{}*{4}{L}@{}}
\optlabel{A}\,$0.18$ & \optlabel{B}\,$0.36$ & \optlabel{C}\,$0.10$ & \optlabel{D}\,$0.20$
\end{tabularx}
\end{qitem}
\qsep

\needspace{7\baselineskip}
\begin{qitem}{26}
A qubit with $T_1 = 50\,\mu$s starts in $\ket1$. After what time is $P(1) = \frac12$?
\par\smallskip\noindent\begin{tabularx}{\linewidth}{@{}*{4}{L}@{}}
\optlabel{A}\,$50\,\mu$s & \optlabel{B}\,$69.3\,\mu$s & \optlabel{C}\,$25\,\mu$s & \optlabel{D}\,$34.7\,\mu$s
\end{tabularx}
\end{qitem}
\qsep

\needspace{7\baselineskip}
\begin{qitem}{27}
A qubit has $T_1 = 100\,\mu$s and pure-dephasing time $T_\varphi = 200\,\mu$s. Using $\frac{1}{T_2} = \frac{1}{2T_1} + \frac{1}{T_\varphi}$, $T_2$ equals
\par\smallskip\noindent\begin{tabularx}{\linewidth}{@{}*{4}{L}@{}}
\optlabel{A}\,$150\,\mu$s & \optlabel{B}\,$66.7\,\mu$s & \optlabel{C}\,$200\,\mu$s & \optlabel{D}\,$100\,\mu$s
\end{tabularx}
\end{qitem}
\qsep

\needspace{7\baselineskip}
\begin{qitem}{28}
Which pair $(T_1, T_2)$ is \emph{physically impossible} in the standard model?
\par\smallskip\noindent\begin{tabularx}{\linewidth}{@{}*{2}{L}@{}}
\optlabel{A}\,$T_1 = 100\,\mu$s, $T_2 = 30\,\mu$s & \optlabel{B}\,$T_1 = 50\,\mu$s, $T_2 = 120\,\mu$s \\[3pt]
\optlabel{C}\,$T_1 = 50\,\mu$s, $T_2 = 100\,\mu$s & \optlabel{D}\,$T_1 = 80\,\mu$s, $T_2 = 80\,\mu$s
\end{tabularx}
\end{qitem}
\qsep

\needspace{7\baselineskip}
\begin{qitem}{29}
Gates take $50$ ns and the qubit has $T_2 = 100\,\mu$s. Requiring the coherence factor $e^{-t/T_2}$ to stay at least $0.9$, the largest number of
sequential gates is about
\par\smallskip\noindent\begin{tabularx}{\linewidth}{@{}*{4}{L}@{}}
\optlabel{A}\,$210$ & \optlabel{B}\,$21$ & \optlabel{C}\,$2000$ & \optlabel{D}\,$105$
\end{tabularx}
\end{qitem}
\qsep

\needspace{7\baselineskip}
\begin{qitem}{30}
A detector reads $1$ when the qubit is in $\ket0$ with probability $0.02$, and reads $0$ when it is in $\ket1$ with probability $0.05$. For a perfectly
prepared $\ket+$, the probability of reading $1$ is
\par\smallskip\noindent\begin{tabularx}{\linewidth}{@{}*{4}{L}@{}}
\optlabel{A}\,$0.515$ & \optlabel{B}\,$0.475$ & \optlabel{C}\,$0.485$ & \optlabel{D}\,$0.500$
\end{tabularx}
\end{qitem}
\qsep

\needspace{7\baselineskip}
\begin{qitem}{31}
A perfect 3-qubit GHZ state is measured with the readout errors of the previous question ($P(1|0) = 0.02$, $P(0|1) = 0.05$, independent per qubit).
The probability of observing $000$ or $111$ is about
\par\smallskip\noindent\begin{tabularx}{\linewidth}{@{}*{4}{L}@{}}
\optlabel{A}\,$0.899$ & \optlabel{B}\,$0.950$ & \optlabel{C}\,$0.857$ & \optlabel{D}\,$0.941$
\end{tabularx}
\end{qitem}
\qsep

\needspace{7\baselineskip}
\begin{qitem}{32}
Which of the following is a \emph{coherent} (unitary) error rather than decoherence or a measurement error?
\begin{itemize}[leftmargin=2.6em, labelsep=0.4em, itemsep=2pt, topsep=3pt]
\item[\optlabel{A}] the discriminator labelling a $\ket1$ outcome as $0$
\item[\optlabel{B}] a systematic over-rotation of every $X$ gate caused by a miscalibrated pulse amplitude
\item[\optlabel{C}] a random phase that differs from shot to shot and averages to dephasing
\item[\optlabel{D}] energy relaxation of $\ket1$ to $\ket0$ (a $T_1$ process)
\end{itemize}
\end{qitem}
\qsep

\needspace{7\baselineskip}
\begin{qitem}{33}
After a CNOT, CUDA-Q's \texttt{Depolarization2(p)} applies each of the 15 non-identity two-qubit Paulis with probability $p/15$. For $p = 0.015$, the
probability that the outcome of a $Z$-basis measurement of the control qubit is flipped by this error is
\par\smallskip\noindent\begin{tabularx}{\linewidth}{@{}*{4}{L}@{}}
\optlabel{A}\,$0.010$ & \optlabel{B}\,$0.008$ & \optlabel{C}\,$0.0075$ & \optlabel{D}\,$0.015$
\end{tabularx}
\end{qitem}
\qsep

\needspace{7\baselineskip}
\begin{qitem}{34}
A 15-qubit GHZ state suffers only phase-flip noise. Which measurement reveals the errors?
\begin{itemize}[leftmargin=2.6em, labelsep=0.4em, itemsep=2pt, topsep=3pt]
\item[\optlabel{A}] the parity $\langle X^{\otimes15}\rangle$ measured after a Hadamard on every qubit
\item[\optlabel{B}] the probability $P(0\cdots0)$ alone
\item[\optlabel{C}] the number of distinct outcomes in the $Z$ basis
\item[\optlabel{D}] the $Z$-basis population $P(0\cdots0) + P(1\cdots1)$
\end{itemize}
\end{qitem}
\qsep

\needspace{7\baselineskip}
\begin{qitem}{35}
Match each channel (List-I) with a set of Kraus operators (List-II).
\par\smallskip\noindent{\renewcommand{\arraystretch}{1.35}%
\begin{tabularx}{\linewidth}{@{}c>{\raggedright\arraybackslash}p{0.33\linewidth}@{\hspace{1.2em}}cL@{}}
\toprule
\multicolumn{2}{@{}l}{\textbf{List-I}} & \multicolumn{2}{l@{}}{\textbf{List-II}} \\ \midrule
\textbf{(P)} & bit flip & \textbf{(1)} & $\operatorname{diag}(1, \sqrt{1-\lambda})$, $\operatorname{diag}(0, \sqrt\lambda)$ \\
\textbf{(Q)} & phase flip & \textbf{(2)} & $\sqrt{1-p}\,I$, $\sqrt p\,Z$ \\
\textbf{(R)} & amplitude damping & \textbf{(3)} & $\begin{psmallmatrix}1&0\\0&\sqrt{1-\gamma}\end{psmallmatrix}$, $\begin{psmallmatrix}0&\sqrt\gamma\\0&0\end{psmallmatrix}$ \\
\textbf{(S)} & phase damping & \textbf{(4)} & $\sqrt{1-p}\,I$, $\sqrt p\,X$ \\
\bottomrule
\end{tabularx}}
\par\smallskip\noindent\begin{tabularx}{\linewidth}{@{}*{2}{L}@{}}
\optlabel{A}\,P$\to$4, Q$\to$1, R$\to$3, S$\to$2 & \optlabel{B}\,P$\to$4, Q$\to$2, R$\to$3, S$\to$1 \\[3pt]
\optlabel{C}\,P$\to$2, Q$\to$4, R$\to$3, S$\to$1 & \optlabel{D}\,P$\to$4, Q$\to$2, R$\to$1, S$\to$3
\end{tabularx}
\end{qitem}
\qsep

\needspace{7\baselineskip}
\begin{qitem}{36}
Match each channel (List-I) with its action on the Bloch vector $(r_x, r_y, r_z)$ (List-II).
\par\smallskip\noindent{\renewcommand{\arraystretch}{1.35}%
\begin{tabularx}{\linewidth}{@{}c>{\raggedright\arraybackslash}p{0.33\linewidth}@{\hspace{1.2em}}cL@{}}
\toprule
\multicolumn{2}{@{}l}{\textbf{List-I}} & \multicolumn{2}{l@{}}{\textbf{List-II}} \\ \midrule
\textbf{(P)} & depolarizing (CUDA-Q convention, $p$) & \textbf{(1)} & $r_x, r_y$ scaled, $r_z$ unchanged \\
\textbf{(Q)} & amplitude damping ($\gamma$) & \textbf{(2)} & all three components multiplied by the same factor $1 - \frac{4p}{3}$ \\
\textbf{(R)} & phase damping ($\lambda$) & \textbf{(3)} & $r_x$ unchanged, $r_y$ and $r_z$ multiplied by $1 - 2p$ \\
\textbf{(S)} & bit flip ($p$) & \textbf{(4)} & $r_x, r_y$ multiplied by $\sqrt{1-\gamma}$, $r_z \mapsto \gamma + (1-\gamma)r_z$ \\
\bottomrule
\end{tabularx}}
\par\smallskip\noindent\begin{tabularx}{\linewidth}{@{}*{2}{L}@{}}
\optlabel{A}\,P$\to$2, Q$\to$1, R$\to$4, S$\to$3 & \optlabel{B}\,P$\to$2, Q$\to$4, R$\to$3, S$\to$1 \\[3pt]
\optlabel{C}\,P$\to$3, Q$\to$4, R$\to$1, S$\to$2 & \optlabel{D}\,P$\to$2, Q$\to$4, R$\to$1, S$\to$3
\end{tabularx}
\end{qitem}
\qsep

\needspace{7\baselineskip}
\begin{qitem}{37}
\textbf{Assertion (A):} Amplitude damping is not unital: it does not map the maximally mixed state $I/2$ to itself.\\[1mm]
\textbf{Reason (R):} Amplitude damping multiplies the off-diagonal elements of $\rho$ by $\sqrt{1 - \gamma}$.
\begin{itemize}[leftmargin=2.6em, labelsep=0.4em, itemsep=2pt, topsep=3pt]
\item[\optlabel{A}] Both Assertion (A) and Reason (R) are true, and R is the correct explanation of A
\item[\optlabel{B}] Both A and R are true, but R is \emph{not} the correct explanation of A
\item[\optlabel{C}] A is true, but R is false
\item[\optlabel{D}] A is false, but R is true
\end{itemize}
\end{qitem}
\qsep

\needspace{7\baselineskip}
\begin{qitem}{38}
\textbf{Statement I:} The phase-flip channel with $p = \frac12$ maps every single-qubit state to a state with zero off-diagonal elements.\\[1mm]
\textbf{Statement II:} The phase-flip channel with $p = 1$ destroys all quantum information.
\begin{itemize}[leftmargin=2.6em, labelsep=0.4em, itemsep=2pt, topsep=3pt]
\item[\optlabel{A}] Both Statement I and Statement II are true
\item[\optlabel{B}] Both Statement I and Statement II are false
\item[\optlabel{C}] Statement I is true, but Statement II is false
\item[\optlabel{D}] Statement I is false, but Statement II is true
\end{itemize}
\end{qitem}
\qsep

\needspace{20\baselineskip}
\begin{passagebox}{Paragraph for Questions 39 and 40}
A superconducting qubit has $T_1 = 80\,\mu$s and $T_2 = 60\,\mu$s. It is prepared in $\ket+$ and then idles for $t = 20\,\mu$s.
Use the standard model: populations relax towards $\ket0$ as $\rho_{11}(t) = \rho_{11}(0)e^{-t/T_1}$, and the coherence decays as
$|\rho_{01}(t)| = |\rho_{01}(0)|e^{-t/T_2}$, with $\frac{1}{T_2} = \frac{1}{2T_1} + \frac{1}{T_\varphi}$ ($T_\varphi$ = pure-dephasing time).
\end{passagebox}
\needspace{7\baselineskip}
\begin{qitem}{39}
After the idle period, $\rho_{11}$ and $|\rho_{01}|$ are approximately
\par\smallskip\noindent\begin{tabularx}{\linewidth}{@{}*{4}{L}@{}}
\optlabel{A}\,$0.389$ and $0.358$ & \optlabel{B}\,$0.500$ and $0.358$ & \optlabel{C}\,$0.389$ and $0.316$ & \optlabel{D}\,$0.389$ and $0.389$
\end{tabularx}
\end{qitem}
\qsep

\needspace{7\baselineskip}
\begin{qitem}{40}
The fidelity $\braket{+|\rho|+}$ after the idle period, and the pure-dephasing time $T_\varphi$, are
\par\smallskip\noindent\begin{tabularx}{\linewidth}{@{}*{2}{L}@{}}
\optlabel{A}\,$0.889$ and $96\,\mu$s & \optlabel{B}\,$0.716$ and $96\,\mu$s \\[3pt]
\optlabel{C}\,$0.858$ and $96\,\mu$s & \optlabel{D}\,$0.858$ and $43.6\,\mu$s
\end{tabularx}
\end{qitem}
\qsep

\newpage

\begin{sectionbanner}SECTION C \enspace$\cdot$\enspace Noise Models and Simulation \hfill Q.41--60\end{sectionbanner}

{\small \textbf{Syllabus:} Kraus operators, quantum channels and CPTP maps, density-matrix evolution, noisy circuits, sampling, gate noise, simulator targets and memory, figures of merit, and the code and recorded results of \texttt{Noise\_Simulation.ipynb}.\\ \textbf{Formats:} Single-correct MCQ (Q.41--54); Matching lists (Q.55--56); Assertion--Reason (Q.57); Statement I / II (Q.58); Paragraph based (Q.59--60). Each question: $+3$ / $0$ / $-1$.}
\par\medskip

\needspace{7\baselineskip}
\begin{qitem}{41}
On the target \texttt{density-matrix-cpu} a kernel prepares $\ket0$ and applies \texttt{rx(t, q[0])}, with \texttt{t = 0.0} passed as a kernel
argument. The noise model attaches \texttt{cudaq.DepolarizationChannel(0.3)} to every \texttt{rx}. \texttt{cudaq.observe} of $Z$ returns
\par\smallskip\noindent\begin{tabularx}{\linewidth}{@{}*{4}{L}@{}}
\optlabel{A}\,$0.7$ & \optlabel{B}\,$0.6$ & \optlabel{C}\,$0.4$ & \optlabel{D}\,$1.0$
\end{tabularx}
\end{qitem}
\qsep

\needspace{7\baselineskip}
\begin{qitem}{42}
The same experiment is repeated, but the kernel now contains \texttt{rx(0.0, q[0])} with the angle written as a \emph{literal constant}. With CUDA-Q 0.16
the result is $\langle Z\rangle = 1.0$. The best explanation is
\begin{itemize}[leftmargin=2.6em, labelsep=0.4em, itemsep=2pt, topsep=3pt]
\item[\optlabel{A}] the compiler removes the constant identity rotation, so there is no \texttt{rx} gate left to carry the noise channel
\item[\optlabel{B}] noise channels are applied only before a gate, never after it
\item[\optlabel{C}] \texttt{density-matrix-cpu} ignores noise models
\item[\optlabel{D}] the depolarizing channel acts only on non-zero rotation angles by definition
\end{itemize}
\end{qitem}
\qsep

\needspace{7\baselineskip}
\begin{qitem}{43}
Which set of operators defines a valid (completely positive, trace-preserving) quantum channel?
\begin{itemize}[leftmargin=2.6em, labelsep=0.4em, itemsep=2pt, topsep=3pt]
\item[\optlabel{A}] $\Big\{\begin{psmallmatrix}1&0\\0&0.6\end{psmallmatrix},\ \begin{psmallmatrix}0&0.6\\0&0\end{psmallmatrix}\Big\}$
\item[\optlabel{B}] $\{I,\ \sqrt{0.1}\,X\}$
\item[\optlabel{C}] $\{\sqrt{0.5}\,I,\ \sqrt{0.5}\,X,\ \sqrt{0.5}\,Z\}$
\item[\optlabel{D}] $\{\sqrt{0.7}\,I,\ \sqrt{0.3}\,Y\}$
\end{itemize}
\end{qitem}
\qsep

\needspace{7\baselineskip}
\begin{qitem}{44}
The bit-flip channel with $p = 0.2$ acts on $\rho = \begin{pmatrix}0.9 & 0.3\\ 0.3 & 0.1\end{pmatrix}$. The output is
\par\smallskip\noindent\begin{tabularx}{\linewidth}{@{}*{4}{L}@{}}
\optlabel{A}\,$\begin{pmatrix}0.74 & 0.18\\ 0.18 & 0.26\end{pmatrix}$ & \optlabel{B}\,$\begin{pmatrix}0.74 & 0.30\\ 0.30 & 0.26\end{pmatrix}$ & \optlabel{C}\,$\begin{pmatrix}0.82 & 0.30\\ 0.30 & 0.18\end{pmatrix}$ & \optlabel{D}\,$\begin{pmatrix}0.90 & 0.18\\ 0.18 & 0.10\end{pmatrix}$
\end{tabularx}
\end{qitem}
\qsep

\needspace{7\baselineskip}
\begin{qitem}{45}
Two depolarizing channels (CUDA-Q convention) with $p_1 = 0.03$ and $p_2 = 0.06$ act one after the other. They equal a single depolarizing channel with $p$ equal to
\par\smallskip\noindent\begin{tabularx}{\linewidth}{@{}*{4}{L}@{}}
\optlabel{A}\,$0.1168$ & \optlabel{B}\,$0.0882$ & \optlabel{C}\,$0.0876$ & \optlabel{D}\,$0.0900$
\end{tabularx}
\end{qitem}
\qsep

\needspace{7\baselineskip}
\begin{qitem}{46}
In \texttt{Noise\_Simulation.ipynb} the 15-qubit GHZ circuit has $G$ noisy gate locations (1 Hadamard and 14 CNOTs, each followed by a depolarizing channel).
The notebook's no-error probability $(1 - p)^G$ at $p = 0.02$ is
\par\smallskip\noindent\begin{tabularx}{\linewidth}{@{}*{4}{L}@{}}
\optlabel{A}\,$0.7000$ & \optlabel{B}\,$0.9800$ & \optlabel{C}\,$0.7386$ & \optlabel{D}\,$0.5455$
\end{tabularx}
\end{qitem}
\qsep

\needspace{7\baselineskip}
\begin{qitem}{47}
For \texttt{NOISE\_MODEL = "bit\_flip"}, the notebook's CNOT channel is \texttt{Pauli2} with probabilities $p(1-p)$ for $IX$ and $XI$ and $p^2$ for $XX$. The
no-error probability of the 15-qubit GHZ circuit at $p = 0.01$ is
\par\smallskip\noindent\begin{tabularx}{\linewidth}{@{}*{4}{L}@{}}
\optlabel{A}\,$0.7472$ & \optlabel{B}\,$0.8601$ & \optlabel{C}\,$0.7397$ & \optlabel{D}\,$0.9801$
\end{tabularx}
\end{qitem}
\qsep

\needspace{7\baselineskip}
\begin{qitem}{48}
The density-matrix method stores $4^N$ complex numbers of 16 bytes. The memory for $N = 13$ qubits is about
\par\smallskip\noindent\begin{tabularx}{\linewidth}{@{}*{4}{L}@{}}
\optlabel{A}\,$1.07$ GB & \optlabel{B}\,$131$ kB & \optlabel{C}\,$268$ MB & \optlabel{D}\,$4.29$ GB
\end{tabularx}
\end{qitem}
\qsep

\needspace{7\baselineskip}
\begin{qitem}{49}
Which simulation can the stabilizer target \texttt{stim} perform \emph{correctly}?
\begin{itemize}[leftmargin=2.6em, labelsep=0.4em, itemsep=2pt, topsep=3pt]
\item[\optlabel{A}] a GHZ circuit with amplitude-damping noise
\item[\optlabel{B}] a GHZ circuit followed by a $T$ gate, with bit-flip noise
\item[\optlabel{C}] a GHZ circuit ($H$, CNOT) with depolarizing noise after every gate
\item[\optlabel{D}] the notebook's mirror circuit with random $R_Y$, $R_Z$ rotations and depolarizing noise
\end{itemize}
\end{qitem}
\qsep

\needspace{7\baselineskip}
\begin{qitem}{50}
On the density-matrix path the notebook calls \texttt{as\_counts\_from\_rho(rho, n, shots)}. What does this function do?
\begin{itemize}[leftmargin=2.6em, labelsep=0.4em, itemsep=2pt, topsep=3pt]
\item[\optlabel{A}] it takes the diagonal of $\rho$ and returns it unchanged as counts, without sampling
\item[\optlabel{B}] it takes the diagonal of the exact noisy $\rho$ as outcome probabilities, reorders the indices so that qubit 0 is written first, and draws \texttt{shots} counts from a multinomial distribution
\item[\optlabel{C}] it measures the density-matrix simulator \texttt{shots} times with \texttt{cudaq.sample}
\item[\optlabel{D}] it rounds every entry of $\rho$ to the nearest multiple of $1/\texttt{shots}$
\end{itemize}
\end{qitem}
\qsep

\needspace{7\baselineskip}
\begin{qitem}{51}
The ideal distribution of a 2-qubit circuit is $P(00) = P(11) = 0.5$. The noisy distribution is $Q(00) = Q(11) = 0.45$, $Q(01) = Q(10) = 0.05$. The Hellinger
fidelity $F_H = \big(\sum_x\sqrt{P_xQ_x}\big)^2$ and the total variation distance $\frac12\sum_x|P_x - Q_x|$ are
\par\smallskip\noindent\begin{tabularx}{\linewidth}{@{}*{4}{L}@{}}
\optlabel{A}\,$0.900$ and $0.200$ & \optlabel{B}\,$0.949$ and $0.100$ & \optlabel{C}\,$0.810$ and $0.100$ & \optlabel{D}\,$0.900$ and $0.100$
\end{tabularx}
\end{qitem}
\qsep

\needspace{7\baselineskip}
\begin{qitem}{52}
A 2-qubit GHZ state is measured after Hadamards on both qubits ($X$ basis): counts \texttt{\{00: 450, 11: 400, 01: 80, 10: 70\}}. The parity
$\langle X\otimes X\rangle$ is
\par\smallskip\noindent\begin{tabularx}{\linewidth}{@{}*{4}{L}@{}}
\optlabel{A}\,$0.40$ & \optlabel{B}\,$0.15$ & \optlabel{C}\,$0.70$ & \optlabel{D}\,$0.85$
\end{tabularx}
\end{qitem}
\qsep

\needspace{7\baselineskip}
\begin{qitem}{53}
For $N = 20$ qubits, by what factor does the memory of exact density-matrix simulation exceed that of one trajectory (state vector) in the same precision?
\par\smallskip\noindent\begin{tabularx}{\linewidth}{@{}*{4}{L}@{}}
\optlabel{A}\,$20$ & \optlabel{B}\,about $10^{12}$ & \optlabel{C}\,$2$ & \optlabel{D}\,about $10^6$
\end{tabularx}
\end{qitem}
\qsep

\needspace{7\baselineskip}
\begin{qitem}{54}
In \texttt{build\_noise\_model} the line \texttt{model.add\_all\_qubit\_channel("x", two, num\_controls=1)} means
\begin{itemize}[leftmargin=2.6em, labelsep=0.4em, itemsep=2pt, topsep=3pt]
\item[\optlabel{A}] the channel \texttt{two} is applied once, at the end of the circuit, to the first two qubits
\item[\optlabel{B}] the channel \texttt{two} replaces the CNOT gate
\item[\optlabel{C}] after every single-qubit \texttt{x} gate, the channel \texttt{two} acts on that qubit and its neighbour
\item[\optlabel{D}] after every CNOT (an \texttt{x} gate with one control), the two-qubit channel \texttt{two} acts on its control and target qubits
\end{itemize}
\end{qitem}
\qsep

\needspace{7\baselineskip}
\begin{qitem}{55}
Match each CUDA-Q simulation target (List-I) with its property for noisy simulation (List-II).
\par\smallskip\noindent{\renewcommand{\arraystretch}{1.35}%
\begin{tabularx}{\linewidth}{@{}c>{\raggedright\arraybackslash}p{0.33\linewidth}@{\hspace{1.2em}}cL@{}}
\toprule
\multicolumn{2}{@{}l}{\textbf{List-I}} & \multicolumn{2}{l@{}}{\textbf{List-II}} \\ \midrule
\textbf{(P)} & \texttt{stim} & \textbf{(1)} & Monte-Carlo trajectories on a GPU, memory $\propto 2^N$, statistical \\
\textbf{(Q)} & \texttt{density-matrix-cpu} & \textbf{(2)} & noiseless state vector: not offered by the noise notebook (it ignores noise models) \\
\textbf{(R)} & \texttt{nvidia} with a noise model & \textbf{(3)} & Clifford circuits with Pauli noise, polynomial cost, hundreds of qubits \\
\textbf{(S)} & \texttt{qpp-cpu} & \textbf{(4)} & exact $\rho$, any Kraus channel, memory $16\cdot4^N$ bytes \\
\bottomrule
\end{tabularx}}
\par\smallskip\noindent\begin{tabularx}{\linewidth}{@{}*{2}{L}@{}}
\optlabel{A}\,P$\to$3, Q$\to$4, R$\to$1, S$\to$2 & \optlabel{B}\,P$\to$3, Q$\to$4, R$\to$2, S$\to$1 \\[3pt]
\optlabel{C}\,P$\to$4, Q$\to$3, R$\to$1, S$\to$2 & \optlabel{D}\,P$\to$3, Q$\to$1, R$\to$4, S$\to$2
\end{tabularx}
\end{qitem}
\qsep

\needspace{7\baselineskip}
\begin{qitem}{56}
Match each figure of merit of the notebook (List-I) with what it is sensitive to or requires (List-II).
\par\smallskip\noindent{\renewcommand{\arraystretch}{1.35}%
\begin{tabularx}{\linewidth}{@{}c>{\raggedright\arraybackslash}p{0.33\linewidth}@{\hspace{1.2em}}cL@{}}
\toprule
\multicolumn{2}{@{}l}{\textbf{List-I}} & \multicolumn{2}{l@{}}{\textbf{List-II}} \\ \midrule
\textbf{(P)} & $Z$-basis GHZ population & \textbf{(1)} & needs the full density matrix; also the off-diagonal (coherence) elements \\
\textbf{(Q)} & $X$-basis parity $\langle X^{\otimes N}\rangle$ & \textbf{(2)} & bit flips (errors that change computational-basis outcomes) \\
\textbf{(R)} & state fidelity $\braket{\psi|\rho|\psi}$ & \textbf{(3)} & phase errors, which are invisible in the $Z$ basis \\
\textbf{(S)} & purity $\operatorname{Tr}\rho^2$ & \textbf{(4)} & mixedness of the state, whatever the target \\
\bottomrule
\end{tabularx}}
\par\smallskip\noindent\begin{tabularx}{\linewidth}{@{}*{2}{L}@{}}
\optlabel{A}\,P$\to$2, Q$\to$3, R$\to$4, S$\to$1 & \optlabel{B}\,P$\to$2, Q$\to$1, R$\to$3, S$\to$4 \\[3pt]
\optlabel{C}\,P$\to$2, Q$\to$3, R$\to$1, S$\to$4 & \optlabel{D}\,P$\to$3, Q$\to$2, R$\to$1, S$\to$4
\end{tabularx}
\end{qitem}
\qsep

\needspace{7\baselineskip}
\begin{qitem}{57}
\textbf{Assertion (A):} For depolarizing noise, the GHZ state fidelity $\braket{\mathrm{GHZ}|\rho|\mathrm{GHZ}}$ is \emph{exactly} equal to the
no-error probability $(1 - p)^G$.\\[1mm]
\textbf{Reason (R):} For Pauli noise, a run without any error produces the ideal state, and no error occurs with probability $(1 - p)^G$.
\begin{itemize}[leftmargin=2.6em, labelsep=0.4em, itemsep=2pt, topsep=3pt]
\item[\optlabel{A}] Both Assertion (A) and Reason (R) are true, and R is the correct explanation of A
\item[\optlabel{B}] Both A and R are true, but R is \emph{not} the correct explanation of A
\item[\optlabel{C}] A is true, but R is false
\item[\optlabel{D}] A is false, but R is true
\end{itemize}
\end{qitem}
\qsep

\needspace{7\baselineskip}
\begin{qitem}{58}
\textbf{Statement I:} In CUDA-Q 0.16 the density-matrix path of the notebook calls \texttt{set\_noise} before \texttt{get\_state},
because \texttt{get\_state} has no \texttt{noise\_model} argument.\\[1mm]
\textbf{Statement II:} The density-matrix method needs $16\cdot4^N$ bytes, so $N = 14$ needs about $4.3$ GB.
\begin{itemize}[leftmargin=2.6em, labelsep=0.4em, itemsep=2pt, topsep=3pt]
\item[\optlabel{A}] Both Statement I and Statement II are true
\item[\optlabel{B}] Both Statement I and Statement II are false
\item[\optlabel{C}] Statement I is true, but Statement II is false
\item[\optlabel{D}] Statement I is false, but Statement II is true
\end{itemize}
\end{qitem}
\qsep

\needspace{20\baselineskip}
\begin{passagebox}{Paragraph for Questions 59 and 60}
The previously recorded run of \texttt{Noise\_Simulation.ipynb} simulates the 15-qubit GHZ circuit ($H$ on qubit 0 and a chain of
14 CNOTs) with depolarizing noise after every gate on the stabilizer target \texttt{stim} (10\,000 shots, seed 42). Part of its table:
\begin{center}\small
\begin{tabular}{rcccr}\toprule
$p$ & population $P(0\cdots0) + P(1\cdots1)$ & parity $\langle X^{\otimes15}\rangle$ & $(1-p)^G$ & distinct outcomes\\ \midrule
0.01 & 0.8936 & 0.6910 & 0.8601 & 96\\
0.1 & 0.3232 & 0.0160 & 0.2059 & 1165\\ \bottomrule
\end{tabular}
\end{center}
\end{passagebox}
\needspace{7\baselineskip}
\begin{qitem}{59}
The probability mass found \emph{outside} the two ideal GHZ outcomes at $p = 0.01$ and at $p = 0.1$ is
\par\smallskip\noindent\begin{tabularx}{\linewidth}{@{}*{2}{L}@{}}
\optlabel{A}\,$0.1399$ and $0.7941$ & \optlabel{B}\,$0.1064$ and $0.6768$ \\[3pt]
\optlabel{C}\,$0.3090$ and $0.9840$ & \optlabel{D}\,$0.8936$ and $0.3232$
\end{tabularx}
\end{qitem}
\qsep

\needspace{7\baselineskip}
\begin{qitem}{60}
Why is the population at $p = 0.01$ ($0.8936$) \emph{larger} than $(1 - p)^G = 0.8601$, while the parity ($0.6910$) is smaller?
\begin{itemize}[leftmargin=2.6em, labelsep=0.4em, itemsep=2pt, topsep=3pt]
\item[\optlabel{A}] $(1-p)^G$ is an upper bound on the population, so the population must lie below it --- the table contains a mistake
\item[\optlabel{B}] $(1-p)^G$ counts every error, but $Z$-type errors leave the $Z$-basis population intact; the parity is also sensitive to those phase errors (and to Hadamard noise in the $X$-basis circuit)
\item[\optlabel{C}] the stabilizer simulator under-reports errors because it cannot simulate depolarizing noise
\item[\optlabel{D}] 10\,000 shots are too few: with more shots the population would fall to exactly $0.8601$
\end{itemize}
\end{qitem}
\qsep

\newpage

\begin{sectionbanner}SECTION D \enspace$\cdot$\enspace Advanced Noise Analysis and Implementation \hfill Q.61--80\end{sectionbanner}

{\small \textbf{Syllabus:} process and average gate fidelity, error rates and randomized benchmarking, zero-noise extrapolation, readout-error mitigation, dynamical decoupling, quantum error correction, twirling, CUDA-Q implementation and practical limitations.\\ \textbf{Formats:} Single-correct MCQ (Q.61--74); Matching lists (Q.75--76); Assertion--Reason (Q.77); Statement I / II (Q.78); Paragraph based (Q.79--80). Each question: $+3$ / $0$ / $-1$.}
\par\medskip

\needspace{7\baselineskip}
\begin{qitem}{61}
For a single-qubit depolarizing channel (CUDA-Q convention: Kraus operators $\sqrt{1-p}\,I$, $\sqrt{p/3}\,X, Y, Z$) the process (entanglement) fidelity is
$F_{\text{pro}} = 1 - p$. Using $F_{\text{avg}} = \frac{dF_{\text{pro}} + 1}{d + 1}$, the average gate fidelity for $p = 0.003$ is
\par\smallskip\noindent\begin{tabularx}{\linewidth}{@{}*{4}{L}@{}}
\optlabel{A}\,$0.9980$ & \optlabel{B}\,$0.9990$ & \optlabel{C}\,$0.9970$ & \optlabel{D}\,$0.99775$
\end{tabularx}
\end{qitem}
\qsep

\needspace{7\baselineskip}
\begin{qitem}{62}
Single-qubit randomized benchmarking fits the survival probability to $A\alpha^m + B$ ($m$ = sequence length) and finds $\alpha = 0.99$. The error per Clifford
$r = \frac{(d - 1)(1 - \alpha)}{d}$ is
\par\smallskip\noindent\begin{tabularx}{\linewidth}{@{}*{4}{L}@{}}
\optlabel{A}\,$0.010$ & \optlabel{B}\,$0.0067$ & \optlabel{C}\,$0.0025$ & \optlabel{D}\,$0.005$
\end{tabularx}
\end{qitem}
\qsep

\needspace{7\baselineskip}
\begin{qitem}{63}
Zero-noise extrapolation measures $E(1) = 0.80$ at the native noise and $E(2) = 0.64$ at doubled noise. The \emph{linear} extrapolation to $\lambda = 0$ gives
\par\smallskip\noindent\begin{tabularx}{\linewidth}{@{}*{4}{L}@{}}
\optlabel{A}\,$0.88$ & \optlabel{B}\,$1.00$ & \optlabel{C}\,$0.80$ & \optlabel{D}\,$0.96$
\end{tabularx}
\end{qitem}
\qsep

\needspace{7\baselineskip}
\begin{qitem}{64}
ZNE with three noise levels gives $E(1) = 0.80$, $E(2) = 0.65$, $E(3) = 0.53$. The quadratic (Richardson) extrapolation $E(0) = 3E(1) - 3E(2) + E(3)$ equals
\par\smallskip\noindent\begin{tabularx}{\linewidth}{@{}*{4}{L}@{}}
\optlabel{A}\,$0.95$ & \optlabel{B}\,$0.985$ & \optlabel{C}\,$0.98$ & \optlabel{D}\,$0.53$
\end{tabularx}
\end{qitem}
\qsep

\needspace{7\baselineskip}
\begin{qitem}{65}
A circuit has 10 gates of equal error. To amplify noise, 4 of them are ``folded'' ($G \to GG^\dagger G$), the others are left unchanged. Assuming the noise grows with
the number of gates, the effective noise scale factor is
\par\smallskip\noindent\begin{tabularx}{\linewidth}{@{}*{4}{L}@{}}
\optlabel{A}\,$2.2$ & \optlabel{B}\,$1.4$ & \optlabel{C}\,$1.8$ & \optlabel{D}\,$3.0$
\end{tabularx}
\end{qitem}
\qsep

\needspace{7\baselineskip}
\begin{qitem}{66}
Readout calibration of one qubit gives $P(\text{read }0|0) = 0.98$ and $P(\text{read }0|1) = 0.05$. An experiment reads $0$ with frequency $0.60$. The mitigated
probability of the true state $\ket0$ is about
\par\smallskip\noindent\begin{tabularx}{\linewidth}{@{}*{4}{L}@{}}
\optlabel{A}\,$0.604$ & \optlabel{B}\,$0.591$ & \optlabel{C}\,$0.600$ & \optlabel{D}\,$0.612$
\end{tabularx}
\end{qitem}
\qsep

\needspace{7\baselineskip}
\begin{qitem}{67}
Which statement about readout-error mitigation by inverting the calibrated confusion matrix $A$ is correct?
\begin{itemize}[leftmargin=2.6em, labelsep=0.4em, itemsep=2pt, topsep=3pt]
\item[\optlabel{A}] it reduces the number of shots needed, because $A^{-1}$ removes the shot noise
\item[\optlabel{B}] it requires quantum error correction to estimate $A$
\item[\optlabel{C}] it also corrects gate errors that happened before the measurement
\item[\optlabel{D}] $A^{-1}\mathbf q$ can contain small negative ``quasi-probabilities'', and the full $2^N\times2^N$ matrix becomes impractical for many qubits
\end{itemize}
\end{qitem}
\qsep

\needspace{7\baselineskip}
\begin{qitem}{68}
A qubit in $\ket+$ idles for $t$ with a random but \emph{static} detuning $\delta \sim \mathcal N(0, \sigma^2)$, $\sigma = 1$ rad/$\mu$s, so it acquires the phase $\delta t$. At
$t = 2\,\mu$s the average coherence $\langle X\rangle$ without and with a Hahn echo ($X$ pulse at $t/2$ and again at the end) is
\par\smallskip\noindent\begin{tabularx}{\linewidth}{@{}*{4}{L}@{}}
\optlabel{A}\,$0.135$ and $0.135$ & \optlabel{B}\,$0.368$ and $1.0$ & \optlabel{C}\,$0.607$ and $1.0$ & \optlabel{D}\,$0.135$ and $1.0$
\end{tabularx}
\end{qitem}
\qsep

\needspace{7\baselineskip}
\begin{qitem}{69}
Now the dephasing is Markovian: an independent phase-flip channel with $p = 0.02$ acts after each of 20 idle steps on $\ket+$. With a Hahn echo inserted at step 10,
$\langle X\rangle$ at the end is
\par\smallskip\noindent\begin{tabularx}{\linewidth}{@{}*{2}{L}@{}}
\optlabel{A}\,$0$ & \optlabel{B}\,$1.000$ \\[3pt]
\optlabel{C}\,$0.442$, the same as without echo & \optlabel{D}\,$0.668$
\end{tabularx}
\end{qitem}
\qsep

\needspace{7\baselineskip}
\begin{qitem}{70}
In the 3-qubit bit-flip code each qubit flips independently with probability $p$, and encoding and correction are perfect. The logical error probability at $p = 0.1$, and
the value of $p$ above which encoding is worse than an unencoded qubit, are
\par\smallskip\noindent\begin{tabularx}{\linewidth}{@{}*{2}{L}@{}}
\optlabel{A}\,$0.028$ and $p = 0.5$ & \optlabel{B}\,$0.010$ and $p = 0.5$ \\[3pt]
\optlabel{C}\,$0.028$ and $p = 1/3$ & \optlabel{D}\,$0.030$ and $p = 0.5$
\end{tabularx}
\end{qitem}
\qsep

\needspace{7\baselineskip}
\begin{qitem}{71}
A surface-code memory has logical error rate $p_L \approx A(p/p_{\text{th}})^{(d+1)/2}$ for code distance $d$. With $p/p_{\text{th}} = 0.1$, increasing the distance from
$d = 3$ to $d = 5$ reduces $p_L$ by a factor of
\par\smallskip\noindent\begin{tabularx}{\linewidth}{@{}*{2}{L}@{}}
\optlabel{A}\,none: $p_L$ increases by a factor of 10 & \optlabel{B}\,$10$ \\[3pt]
\optlabel{C}\,$2$ & \optlabel{D}\,$100$
\end{tabularx}
\end{qitem}
\qsep

\needspace{7\baselineskip}
\begin{qitem}{72}
Each of $n = 60$ idle gates over-rotates coherently by $\varepsilon = 0.05$ rad about $x$. Compare with a stochastic bit flip of the same per-gate error probability
$p = \sin^2(\varepsilon/2)$. Starting from $\ket0$, the probabilities of finding $1$ after 60 gates are approximately
\par\smallskip\noindent\begin{tabularx}{\linewidth}{@{}*{2}{L}@{}}
\optlabel{A}\,$0.995$ (coherent) and $0.036$ (stochastic) & \optlabel{B}\,$0.0375$ and $0.036$ \\[3pt]
\optlabel{C}\,$0.995$ and $0.995$ & \optlabel{D}\,$0.036$ and $0.036$
\end{tabularx}
\end{qitem}
\qsep

\needspace{7\baselineskip}
\begin{qitem}{73}
Which statement about Pauli twirling (randomized compiling) is correct?
\begin{itemize}[leftmargin=2.6em, labelsep=0.4em, itemsep=2pt, topsep=3pt]
\item[\optlabel{A}] it removes all errors, so the circuit becomes noiseless
\item[\optlabel{B}] averaging a gate over random Pauli frames that cancel ideally turns its coherent error into a stochastic Pauli channel with the same average fidelity
\item[\optlabel{C}] it requires a density-matrix simulator and cannot be done on hardware
\item[\optlabel{D}] it converts stochastic Pauli noise into coherent errors to make them easier to correct
\end{itemize}
\end{qitem}
\qsep

\needspace{7\baselineskip}
\begin{qitem}{74}
The noise notebook runs on a CPU (no GPU visible) with the settings\\
\hspace*{1em}\texttt{CIRCUIT = "ghz"},\ \ \texttt{NOISE\_MODEL = "amplitude\_damping"},\ \ \texttt{SIMULATION\_METHOD = "auto"}.\\
Which method does it choose?
\begin{itemize}[leftmargin=2.6em, labelsep=0.4em, itemsep=2pt, topsep=3pt]
\item[\optlabel{A}] \texttt{stabilizer} (target \texttt{stim})
\item[\optlabel{B}] \texttt{trajectory} (target \texttt{nvidia})
\item[\optlabel{C}] \texttt{statevector} (target \texttt{qpp-cpu})
\item[\optlabel{D}] \texttt{density\_matrix} (target \texttt{density-matrix-cpu})
\end{itemize}
\end{qitem}
\qsep

\needspace{7\baselineskip}
\begin{qitem}{75}
Match each technique (List-I) with its description (List-II).
\par\smallskip\noindent{\renewcommand{\arraystretch}{1.35}%
\begin{tabularx}{\linewidth}{@{}c>{\raggedright\arraybackslash}p{0.33\linewidth}@{\hspace{1.2em}}cL@{}}
\toprule
\multicolumn{2}{@{}l}{\textbf{List-I}} & \multicolumn{2}{l@{}}{\textbf{List-II}} \\ \midrule
\textbf{(P)} & zero-noise extrapolation & \textbf{(1)} & invert a calibrated confusion matrix applied to the measured distribution \\
\textbf{(Q)} & readout-error mitigation & \textbf{(2)} & run the circuit at amplified noise levels and extrapolate the observable to zero noise \\
\textbf{(R)} & dynamical decoupling & \textbf{(3)} & encode logical qubits redundantly, measure syndromes and correct; needs physical errors below a threshold \\
\textbf{(S)} & quantum error correction & \textbf{(4)} & apply pulse sequences during idle periods to refocus slow (quasi-static) dephasing \\
\bottomrule
\end{tabularx}}
\par\smallskip\noindent\begin{tabularx}{\linewidth}{@{}*{2}{L}@{}}
\optlabel{A}\,P$\to$2, Q$\to$1, R$\to$4, S$\to$3 & \optlabel{B}\,P$\to$1, Q$\to$2, R$\to$4, S$\to$3 \\[3pt]
\optlabel{C}\,P$\to$2, Q$\to$1, R$\to$3, S$\to$4 & \optlabel{D}\,P$\to$2, Q$\to$4, R$\to$1, S$\to$3
\end{tabularx}
\end{qitem}
\qsep

\needspace{7\baselineskip}
\begin{qitem}{76}
Match each physical cause (List-I) with its category (List-II).
\par\smallskip\noindent{\renewcommand{\arraystretch}{1.35}%
\begin{tabularx}{\linewidth}{@{}c>{\raggedright\arraybackslash}p{0.33\linewidth}@{\hspace{1.2em}}cL@{}}
\toprule
\multicolumn{2}{@{}l}{\textbf{List-I}} & \multicolumn{2}{l@{}}{\textbf{List-II}} \\ \midrule
\textbf{(P)} & a pulse amplitude 1\,\% too large on every $X$ gate & \textbf{(1)} & decoherence (non-unital, $T_1$) \\
\textbf{(Q)} & a random Pauli error after each gate & \textbf{(2)} & classical measurement (readout) error \\
\textbf{(R)} & energy loss of $\ket1$ to the environment & \textbf{(3)} & coherent error (unitary, systematic) \\
\textbf{(S)} & the discriminator assigns some $\ket1$ shots to outcome 0 & \textbf{(4)} & stochastic (incoherent) gate noise \\
\bottomrule
\end{tabularx}}
\par\smallskip\noindent\begin{tabularx}{\linewidth}{@{}*{2}{L}@{}}
\optlabel{A}\,P$\to$4, Q$\to$3, R$\to$1, S$\to$2 & \optlabel{B}\,P$\to$3, Q$\to$4, R$\to$1, S$\to$2 \\[3pt]
\optlabel{C}\,P$\to$3, Q$\to$4, R$\to$2, S$\to$1 & \optlabel{D}\,P$\to$3, Q$\to$1, R$\to$4, S$\to$2
\end{tabularx}
\end{qitem}
\qsep

\needspace{7\baselineskip}
\begin{qitem}{77}
\textbf{Assertion (A):} Zero-noise extrapolation increases the number of circuit executions needed to reach a given statistical precision.\\[1mm]
\textbf{Reason (R):} ZNE corrects errors on the quantum state during the computation, in the same way as quantum error correction.
\begin{itemize}[leftmargin=2.6em, labelsep=0.4em, itemsep=2pt, topsep=3pt]
\item[\optlabel{A}] Both Assertion (A) and Reason (R) are true, and R is the correct explanation of A
\item[\optlabel{B}] Both A and R are true, but R is \emph{not} the correct explanation of A
\item[\optlabel{C}] A is true, but R is false
\item[\optlabel{D}] A is false, but R is true
\end{itemize}
\end{qitem}
\qsep

\needspace{7\baselineskip}
\begin{qitem}{78}
\textbf{Statement I:} \texttt{Noise\_Simulation.ipynb} includes measurement (readout) errors in its noise model.\\[1mm]
\textbf{Statement II:} The recorded results of \texttt{Noise\_Simulation.ipynb} were obtained on a physical quantum processor.
\begin{itemize}[leftmargin=2.6em, labelsep=0.4em, itemsep=2pt, topsep=3pt]
\item[\optlabel{A}] Both Statement I and Statement II are true
\item[\optlabel{B}] Both Statement I and Statement II are false
\item[\optlabel{C}] Statement I is true, but Statement II is false
\item[\optlabel{D}] Statement I is false, but Statement II is true
\end{itemize}
\end{qitem}
\qsep

\needspace{20\baselineskip}
\begin{passagebox}{Paragraph for Questions 79 and 80}
For the lecture (illustrative experiment, not part of the notebook), zero-noise extrapolation (ZNE) was applied to the parity
$\langle X^{\otimes4}\rangle$ of a 4-qubit GHZ state. The circuit was simulated exactly on \texttt{density-matrix-cpu} with the notebook's
depolarizing noise model at the scaled probability $p = 0.01\lambda$ (noise scale factor $\lambda$). The ideal value is $1$.
\begin{center}\small
\begin{tabular}{cccc}\toprule $\lambda$ & 1 & 2 & 3\\ \midrule $\langle X^{\otimes4}\rangle$ & 0.9055 & 0.8189 & 0.7396\\ \bottomrule\end{tabular}
\end{center}
\end{passagebox}
\needspace{7\baselineskip}
\begin{qitem}{79}
The linear extrapolation from $\lambda = 1, 2$ and the Richardson (quadratic) extrapolation from $\lambda = 1, 2, 3$ give
\par\smallskip\noindent\begin{tabularx}{\linewidth}{@{}*{2}{L}@{}}
\optlabel{A}\,$0.9921$ and $0.9994$ & \optlabel{B}\,$0.9055$ and $0.9921$ \\[3pt]
\optlabel{C}\,$0.9921$ and $1.0021$ & \optlabel{D}\,$0.9872$ and $0.9994$
\end{tabularx}
\end{qitem}
\qsep

\needspace{7\baselineskip}
\begin{qitem}{80}
Which conclusion about this ZNE experiment is correct?
\begin{itemize}[leftmargin=2.6em, labelsep=0.4em, itemsep=2pt, topsep=3pt]
\item[\optlabel{A}] on hardware the same result is obtained without scaling the noise, since $\lambda$ is a software parameter
\item[\optlabel{B}] ZNE has corrected the GHZ state itself, so a subsequent circuit could use the improved state
\item[\optlabel{C}] the extrapolated values depend on the fitted model (an exponential fit gives $1.0021 > 1$, unphysical) and, with shot noise, would carry larger error bars than the raw value
\item[\optlabel{D}] because the exact density matrix was used, ZNE is no longer needed: the raw value $0.9055$ is already the ideal one
\end{itemize}
\end{qitem}
\qsep

\newpage

\end{document}